% GENERATED FILE. Edit canonical lessons, then run npm run generate:books.
% canonical-source-hash: fnv1a64:4fc7bbdfd57f1e9f
% canonical-lessons: JA-C93-machidooshii, JA-C93-inori, JA-C93-chuumon, JA-C93-kouho, JA-C93-yotei

\chapter{Long-awaited, Prayer, Order, Candidate, Plan}
\label{ch:ja-long-awaited-prayer-order-candidate-plan}
\hlchaptermodality{93}

\begin{chapteropening}
\textbf{By the end of this chapter:} \emph{I can say long-awaited, a prayer, an order, a candidate and a plan.}

Complete the last lesson of chapter 93: i can say long-awaited, a prayer, an order, a candidate and a plan.
\end{chapteropening}

\section[\texorpdfstring{\ja{まちどおしい} (machidooshii)}{まちどおしい (machidooshii)}]{\ja{まちどおしい} (machidooshii) — long-awaited}

\label{lesson:JA-C93-machidooshii}



\begin{quote}
Before the new one: say the Japanese for a request, then the Japanese for a desire.
\end{quote}

\subsection*{You'll want to know: \ja{まちどおしい}}
\textbf{\ja{まちどおしい}} — \emph{machidooshii} — \textquotedblleft{}long-awaited\textquotedblright{}.

\begin{grammarlens}[title={What you've built}]
One more word for this chapter.
\end{grammarlens}

\subsection*{Your turn}
\begin{itemize}
\raggedright
  \item \emph{Say it:} \emph{machidooshii}
  \item \emph{Say it:} \emph{machidooshii}, once more
  \item \emph{Say it:} say \emph{machidooshii}
  \item \emph{Recall:} say the Japanese for a request, then the Japanese for a desire, then say \emph{machidooshii} again
\end{itemize}

\begin{tcolorbox}[breakable,colback=teal!4,colframe=teal!35!black,title={Before you move on}]
What does \ja{まちどおしい} mean? (\textquotedblleft{}Long-awaited\textquotedblright{}.) Say it once more.
\end{tcolorbox}
\section[\texorpdfstring{\ja{いのり} (inori)}{いのり (inori)}]{\ja{いのり} (inori) — a prayer}

\label{lesson:JA-C93-inori}



\begin{quote}
Before the new one: say the Japanese for a desire, then the Japanese for long-awaited.
\end{quote}

\subsection*{You'll want to know: \ja{いのり}}
\textbf{\ja{いのり}} — \emph{inori} — \textquotedblleft{}a prayer\textquotedblright{}.

\begin{grammarlens}[title={What you've built}]
One more word for this chapter.
\end{grammarlens}

\subsection*{Your turn}
\begin{itemize}
\raggedright
  \item \emph{Say it:} \emph{inori}
  \item \emph{Say it:} \emph{inori}, once more
  \item \emph{Say it:} say \emph{inori}
  \item \emph{Recall:} say the Japanese for a desire, then the Japanese for long-awaited, then say \emph{inori} again
\end{itemize}

\begin{tcolorbox}[breakable,colback=teal!4,colframe=teal!35!black,title={Before you move on}]
What does \ja{いのり} mean? (\textquotedblleft{}A prayer\textquotedblright{}.) Say it once more.
\end{tcolorbox}
\section[\texorpdfstring{\ja{ちゅうもん} (chuumon)}{ちゅうもん (chuumon)}]{\ja{ちゅうもん} (chuumon) — an order}

\label{lesson:JA-C93-chuumon}



\begin{quote}
Before the new one: say the Japanese for long-awaited, then the Japanese for a prayer.
\end{quote}

\subsection*{You'll want to know: \ja{ちゅうもん}}
\textbf{\ja{ちゅうもん}} — \emph{chuumon} — \textquotedblleft{}an order\textquotedblright{}.

\begin{grammarlens}[title={What you've built}]
One more word for this chapter.
\end{grammarlens}

\subsection*{Your turn}
\begin{itemize}
\raggedright
  \item \emph{Say it:} \emph{chuumon}
  \item \emph{Say it:} \emph{chuumon}, once more
  \item \emph{Say it:} say \emph{chuumon}
  \item \emph{Recall:} say the Japanese for long-awaited, then the Japanese for a prayer, then say \emph{chuumon} again
\end{itemize}

\begin{tcolorbox}[breakable,colback=teal!4,colframe=teal!35!black,title={Before you move on}]
What does \ja{ちゅうもん} mean? (\textquotedblleft{}An order\textquotedblright{}.) Say it once more.
\end{tcolorbox}
\section[\texorpdfstring{\ja{こうほ} (kouho)}{こうほ (kouho)}]{\ja{こうほ} (kouho) — a candidate, a choice}

\label{lesson:JA-C93-kouho}



\begin{quote}
Before the new one: say the Japanese for a prayer, then the Japanese for an order.
\end{quote}

\subsection*{You'll want to know: \ja{こうほ}}
\textbf{\ja{こうほ}} — \emph{kouho} — \textquotedblleft{}a candidate, a choice\textquotedblright{}.

\begin{grammarlens}[title={What you've built}]
One more word for this chapter.
\end{grammarlens}

\subsection*{Your turn}
\begin{itemize}
\raggedright
  \item \emph{Say it:} \emph{kouho}
  \item \emph{Say it:} \emph{kouho}, once more
  \item \emph{Say it:} say \emph{kouho}
  \item \emph{Recall:} say the Japanese for a prayer, then the Japanese for an order, then say \emph{kouho} again
\end{itemize}

\begin{tcolorbox}[breakable,colback=teal!4,colframe=teal!35!black,title={Before you move on}]
What does \ja{こうほ} mean? (\textquotedblleft{}A candidate, a choice\textquotedblright{}.) Say it once more.
\end{tcolorbox}
\section[\texorpdfstring{\ja{よてい} (yotei)}{よてい (yotei)}]{\ja{よてい} (yotei) — a plan}

\label{lesson:JA-C93-yotei}



\begin{quote}
Before the new one: say the Japanese for an order, then the Japanese for a candidate.
\end{quote}

\subsection*{You'll want to know: \ja{よてい}}
\textbf{\ja{よてい}} — \emph{yotei} — \textquotedblleft{}a plan\textquotedblright{}.

\begin{grammarlens}[title={What you've built}]
One more word for this chapter.
\end{grammarlens}

\subsection*{Your turn}
\begin{itemize}
\raggedright
  \item \emph{Say it:} \emph{yotei}
  \item \emph{Say it:} \emph{yotei}, once more
  \item \emph{Say it:} say \emph{yotei}
  \item \emph{Recall:} say the Japanese for an order, then the Japanese for a candidate, then say \emph{yotei} again
\end{itemize}

\begin{tcolorbox}[breakable,colback=teal!4,colframe=teal!35!black,title={Before you move on}]
What does \ja{よてい} mean? (\textquotedblleft{}A plan\textquotedblright{}.) Say it once more.
\end{tcolorbox}
